\section{Related literature}
\label{sec:literature}

The emulator sits at the intersection of three literatures: the UK tradition of large-scale macroeconometric policy models, the post-crisis debate over what class of model official institutions should operate, and the replication and open-science movement in economics.

\subsection{The UK macroeconometric modelling tradition}

The model being emulated is a direct descendant of the estimated structural models pioneered by \citet{klein1950} at the Cowles Commission: a system of behavioural equations estimated on national-accounts data and closed by accounting identities, solved numerically to produce forecasts and policy simulations. The Treasury built the first version of its UK model in 1970, and the model passed into joint OBR/Treasury ownership when the OBR was created in 2010, with governance set out in a published memorandum of understanding \citep{obrmou, obr2013model}. Through the 1970s and 1980s the UK sustained an unusually rich ecosystem of competing models of this class---the Treasury, the National Institute of Economic and Social Research (NIESR), the London Business School, the Bank of England, and the Liverpool model among them---subjected to systematic comparative evaluation by the ESRC Macroeconomic Modelling Bureau at Warwick \citep{wallis1987}. The Bureau's reviews established a discipline that this paper consciously imitates at small scale: deposit the model, document its properties, decompose its forecasts into model and judgement, and compare simulation responses across models on common shocks. That comparative infrastructure largely disappeared in the 1990s; one reading of the present project is an attempt to rebuild a small piece of it in open source, for the one UK model whose equation listing is still published.

The OBR/Treasury model's closest living relatives are NIESR's NiGEM, a quarterly estimated global model with country blocks linked by trade and financial flows, maintained commercially and used by central banks and finance ministries for scenario analysis \citep{hantzsche2018}, and---across the Atlantic and in a more explicitly optimising idiom---the Federal Reserve's FRB/US \citep{brayton1996}, whose public release of equations, data, and simulation code is the transparency standard this project aims at for the UK. The Bank of England, by contrast, moved its central forecasting platform to an estimated DSGE core (COMPASS) surrounded by a suite of auxiliary models \citep{burgess2013}; the OBR/Treasury model is thus the last large non-DSGE structural econometric model in operational use by a UK macro-fiscal institution.

\subsection{Which model class for policy institutions?}

The choice embodied in the OBR model---error-correction econometrics over microfounded optimisation---was contested territory well before the financial crisis and became more so after it. \citet{wrenlewis2018} argues that the ``microfoundations hegemony'' crowded out precisely the class of structural econometric models (SEMs) to which the OBR model belongs, and that policy analysis would have been better served had SEMs retained equal academic standing. \citet{blanchard2018}, taxonomising models into foundational, DSGE, policy, toy, and forecasting classes, places institutional models of the OBR type in the ``policy'' class, where fitting the data and capturing actual transmission mechanisms trump theoretical purity. \citet{hendry2018} make the econometric case directly: for forecasting and policy in the presence of structural breaks, congruent equilibrium-correction systems---exactly the \citet{davidson1978}/\citet{engle1987} architecture of Section~\ref{sec:ecm}---dominate models whose cross-equation restrictions are derived from misspecified theory. This paper takes no side in that debate; its premise is narrower: whatever the merits of the model class, the model that actually produces the UK's official fiscal forecast should be publicly runnable. The multi-model design of the PolicyEngine Macro suite---macroeconometric, overlapping-generations, and structural-VAR members scoring the same reform---is itself an implementation of the ``suite of models'' philosophy that \citet{burgess2013} describe at the Bank of England.

\subsection{Replication, emulation, and open models}

Finally, the project draws method and motivation from the replication literature. Systematic replication attempts find that even published academic results frequently cannot be reproduced from the authors' own materials \citep{chang2015}, and the best-known failure in the fiscal-policy domain---the spreadsheet error in Reinhart and Rogoff's debt-threshold result, found only when a graduate student obtained and re-ran the original file \citep{herndon2014}---is a standing argument that consequential numbers should be recomputable by outsiders. Official fiscal scoring is a stronger case still: OBR forecasts condition UK budget decisions directly, yet the model that produces them has never been publicly executable. The emulator is best understood as a \emph{replication-as-infrastructure} exercise: not a one-off audit of a published result, but a maintained, continuously tested re-implementation whose agreement with the official output is machine-checked on every commit (Section~\ref{sec:calibration}), in the same spirit in which PolicyEngine re-implements statutory tax--benefit rules as open code \citep{policyengine}. The honest-scorecard discipline of Section~\ref{sec:rawscore}---publishing the free-running failures alongside the anchored successes, and refusing to tune unpublished constants to hit known answers---is the transparency analogue of pre-registration: it separates what the replication demonstrates from what it merely reproduces.
